\textbf{Observation 1: `Asteroid occultation'.}

Beginning of occultation: 23:03:31\\
End of occultation: 23:03:48

\begin{center}
\begin{tabular}{c|c|c|c}
Mid-time of occultation & $\pm$ error & Duration of occultation & $\pm$ error \\
\hline
23:03:39.5 & 0.4 s & 17 s & 0.2 s \\
\end{tabular}
\end{center}

Marking:
\begin{enumerate}
\item Mid-time $t$ of occultation:
  \begin{itemize}
  \item $23{:}03{:}39.3 \le t < 23{:}03{:}39.7$ $\Longrightarrow$ 4 points
  \item $23{:}03{:}39.0 \le t < 23{:}03{:}39.3$ or
    $23{:}03{:}39.7 < t \le 23{:}03{:}40.0$ $\Longrightarrow$ 3 points
  \item $23{:}03{:}38.0 \le t < 23{:}03{:}39.0$ or
    $23{:}03{:}40.0 < t \le 23{:}03{:}41.0$ $\Longrightarrow$ 1 point
  \item outside this range $\Longrightarrow$ 0 points
  \end{itemize}
\item Error in mid-time $\Delta t$:
  \begin{itemize}
  \item $0.1 < \Delta x \le 0.5$ $\Longrightarrow$ 3 points
  \item $0.5 < \Delta x \le 0.7$ $\Longrightarrow$ 2 points
  \item $0.7 < \Delta x \le 1.0$ $\Longrightarrow$ 1 point
  \item outside this range or missing $\Longrightarrow$ 0 points
  \end{itemize}
  % sic: the official solution uses "Delta x" in the bands for the error in
  % mid-time although the heading defines it as Delta t
\item Duration $x$ of occultation (should be to 1 s.f.):
  \begin{itemize}
  \item $16.8~\mathrm{s} \le x \le 17.2~\mathrm{s}$ $\Longrightarrow$ 4 points
  \item $16~\mathrm{s} \le x < 16.8~\mathrm{s}$ or
    $17.2~\mathrm{s} < x \le 18~\mathrm{s}$ $\Longrightarrow$ 3 points
  \item $15~\mathrm{s} \le x < 16~\mathrm{s}$ or
    $18~\mathrm{s} < x \le 19~\mathrm{s}$ $\Longrightarrow$ 1 point
  \item outside this range $\Longrightarrow$ 0 points
  \end{itemize}
\item Error in duration $\Delta x$:
  \begin{itemize}
  \item $0.05 < \Delta x \le 0.2$ $\Longrightarrow$ 3 points
  \item $0.2 < \Delta x \le 0.4$ $\Longrightarrow$ 2 points
  \item $0.4 < \Delta x \le 1.0$ $\Longrightarrow$ 1 point
  \item outside this range or missing $\Longrightarrow$ 0 points
  \end{itemize}
\item Error of duration lower than and different from error in mid time
  ($\Delta x < \Delta t$) $\Longrightarrow$ 1 point
\end{enumerate}

\medskip
\textbf{Observation 2: `Starlink'.}

The satellites are at an altitude $h = 20^\circ$ and their height above the
surface of the Earth is $H = 400$ km.

1) In the (simplified) solution neglecting the Earth's curvature, the
student will assume the distance to the satellites is equal to:
\[
d_{flat} = \frac{H}{\sin h} = 1170~\mathrm{km}
\]

% FIGURE NEEDED: flat-Earth sketch -- right triangle with the observer, the
% satellite at height H and the line of sight d at altitude h
% (2023_TEL_sol.pdf page 5)

2) In the solution taking into account the Earth's curvature, let us draw
the $OSC$ (observer--satellite--center of the Earth) triangle. The angles
are: $\angle COS = 90^\circ + h$, $\angle OSC$ that we shall denote as
$\eta$, and $\angle SCO = 180^\circ - (90^\circ + h) - \eta = 70^\circ - \eta$.

% FIGURE NEEDED: curved-Earth sketch -- the observer O on the spherical
% Earth, satellite S, centre C, with radii R_earth, height H and distance d
% (2023_TEL_sol.pdf page 6)

The law of sines can be applied:
\[
\frac{\sin(90^\circ + h)}{R_\oplus + H} = \frac{\sin\eta}{R_\oplus}
= \frac{\sin(70^\circ - \eta)}{d_{curve}}
\]
We calculate the distance to the satellites:
\[
\eta = \arcsin\left(\frac{R_\oplus}{R_\oplus + H}\sin(90^\circ + h)\right)
= 62^\circ.2
\]
\[
d_{curve} = R_\oplus \frac{\sin(70^\circ - \eta)}{\sin\eta} = 980~\mathrm{km}
\]

Knowing the radius of the orbit is $R_\oplus + H$, we calculate
$v_{orbit} = \sqrt{GM_\oplus/(R_\oplus + H)} = 7.6$ km/s. The observer can
only measure the tangential component of motion $v_t$.

The student will estimate the angular velocity to be
$v_t/d = v_{orbit}\sin h/d = 2.2\cdot 10^{-3}\,[1/\mathrm{s}] = 7.6\,['/\mathrm{s}]$
(for the solution neglecting the Earth's curvature) or
$2.6\cdot 10^{-3}\,[1/\mathrm{s}] = 9.0\,['/\mathrm{s}]$ (for the solution
taking into account the Earth's curvature).

\begin{flushright}
\textbf{5 points for predicting the angular velocity:}\\
5 points if within $8.55 - 9.45\,['/\mathrm{s}]$;\\
4 points if within $8.1 - 9.9\,['/\mathrm{s}]$;\\
3 points if within $7.5 - 10.5\,['/\mathrm{s}]$;\\
0 points otherwise
\end{flushright}

The simulated angular velocity is $8.5\,['/\mathrm{s}]$; the student should
conclude this value is similar to their prediction (or optionally comment it
is slightly lower and might imply e.g.\ the real orbit radius is minimally
larger than predicted, but it is not necessary).

\begin{flushright}
\textbf{4 points for measuring the angular velocity:}\\
4 points if agrees within 5\% ($8.1$ - $8.9$ $['/\mathrm{s}]$);\\
2 points if agrees within 10\% ($7.65$ - $9.35$ $['/\mathrm{s}]$);\\
0 points otherwise\\[4pt]
\textbf{1 point for correct comparison}
\end{flushright}

The student should then measure the time interval between the passes to be
$t = 2$ s.

\begin{flushright}
\textbf{2 points for measuring the time interval:}\\
2 points if agrees within 5\% ($1.9$ - $2.1$ s); 0 points otherwise
\end{flushright}

In such a short time, the path along the orbit can be approximated with a
straight line; the distance passed is $S = v_{orbit}t = 15$ km, which is the
distance between the satellites. (Alternatively, if the student chooses to
use their own $v_t$ measurement, $S = v_t t/\sin h = 14$ km).

\begin{flushright}
\textbf{2 points for calculating the distance}\\
2 points if within 13.8-15.2 km\\
1 point if within 12.5-16.5 km\\
0 points otherwise
\end{flushright}

Finally, the student should mark the observed satellite path on the sky map.

\begin{flushright}
\textbf{1 point for correctly marking the satellite trail}
\end{flushright}

\medskip
\textbf{Observation 3: `Planetary Moons'.}

The displayed planet is Saturn. On the reference solution image the moons
are marked and labelled with their names and magnitudes: Japetus (11.0),
Rhea (9.9), Mimas (13.0), Tethys (10.4), Enceladus (11.8), Dione (10.6) and
Titan (8.5) around the disk of Saturn.

% FIGURE NEEDED: solution image of Saturn with the visible moons circled and
% labelled (Japetus, Rea, Mimas, Tetyda, Enceladus, Dione, Tytan -- labels
% in the original carry slightly different magnitudes: Rea 9.8, Tetyda 10.3,
% Dione 10.5, Enceladus 11.8, Tytan 8.4, Mimas 13, Japetus 11)
% (2023_TEL_sol.pdf page 9)

2 points for each moon: position (1 pt) and name (1 pt).

\medskip
\textbf{Observation 4: `Supernova'.}

\begin{center}
\begin{tabular}{c|c|c}
Right ascension & Declination & est.\ magnitude \\
\hline
$9^h\,55^m\,54^s$ & $+69^\circ\,09'\,11''$ & 10.2 mag \\
\end{tabular}
\end{center}

\begin{enumerate}
\item Declination:
  \begin{itemize}
  \item[(a)] $< \pm 1.5'$ $\Longrightarrow$ 4 points
  \item[(b)] $\le \pm 3'$ $\Longrightarrow$ 2 points
  \item[(c)] $> \pm 3'$ $\Longrightarrow$ 0 points
  \end{itemize}
\item Right ascension:
  \begin{itemize}
  \item[(a)] $< \pm\,22.5$ s $\Longrightarrow$ 4 points
  \item[(b)] $\le \pm\,45$ s $\Longrightarrow$ 2 points
  \item[(c)] $> \pm\,45$ s $\Longrightarrow$ 0 points
  \end{itemize}
\item Magnitude:
  \begin{itemize}
  \item[(a)] $< \pm\,0.4$ mag $\Longrightarrow$ 2 points
  \item[(b)] $\le \pm\,0.8$ mag $\Longrightarrow$ 1 point
  \item[(c)] $> \pm\,0.8$ mag $\Longrightarrow$ 0 points
  \end{itemize}
\end{enumerate}
